%======================================================================
%  Front matter: title page, certificate, declaration, abstract, ToC
%======================================================================
\begin{titlepage}
\centering
\thispagestyle{empty}

{\color{navy}\rule{\textwidth}{1.6pt}}\\[2pt]
{\color{teal}\rule{\textwidth}{0.6pt}}\\[16mm]

{\sffamily\bfseries\LARGE\color{navy}
Channel Aware Authentication Protocol\\[4pt]
for Secure Underwater Communication Systems\\}
\vspace{7mm}
{\sffamily\Large\color{slate} Overview of Previous Work\\[2pt]
and Current Upgrades}\\[16mm]

\end{titlepage}

%----------------------------------------------------------------------


%----------------------------------------------------------------------
\chapter*{Abstract}
\addcontentsline{toc}{chapter}{Abstract}
\thispagestyle{plain}

\noindent
Underwater wireless communication networks carry environmental monitoring, offshore
exploration and naval surveillance traffic, yet they cannot use any of the security
machinery that terrestrial networks rely on. Radio does not propagate in seawater, so these
networks talk in sound: about 1500\,m/s, roughly 10\,kbps of usable bandwidth, 10--30\,\%
packet loss, and battery-powered nodes that nobody can service. A TLS handshake on such a
link would take tens of seconds and a measurable fraction of a node's lifetime energy.

\vspace{2mm}\noindent
This report studies the authentication protocol proposed by Rupa \emph{et al.} in the
\emph{IEEE Internet of Things Journal} (Nov.\ 2025), which is designed specifically for that
regime. The protocol arranges the network as a five-layer chain --- underwater sensor,
submarine, buoy, satellite, base station --- and authenticates \emph{each hop
independently} using elliptic-curve public-key encryption, a fresh nonce and a bounded
timestamp window, completing the seabed-to-shore chain in twelve messages and 2112 bits.
It adds dynamic fallback peers so that a failed buoy does not tear down the session, and it
is validated with BAN logic and the Scyther model checker.

\vspace{2mm}\noindent
We reproduce the protocol as an eight-node Python simulator, then strengthen it in seven
places where our own first version was unrealistic or unsafe. The most consequential change
is cryptographic: our initial build used a hand-written repeating-key XOR cipher, which is
broken by a known-plaintext or key-reuse attack in seconds; we replaced it with
\textbf{AES-256-GCM} authenticated encryption, giving both confidentiality and a 128-bit
integrity tag that makes ciphertext tampering detectable. The remaining six enhancements
replace invented numbers with physics: Thorp acoustic absorption for delay, a 15\,\%
Bernoulli loss channel with retransmission, per-node battery depletion, random-waypoint
mobility for AUVs, a $z$-score detector for delay-injection attacks, and an extended set of
result plots benchmarked against five prior schemes.

\vspace{2mm}\noindent
Separately, we rewrote the Scyther specification. Our first model expressed the chain as one
monolithic five-role protocol and failed several agreement claims; restructuring it into
four independent Needham--Schroeder--Lowe hops, with message tags and the responder identity
inside the second message, produced \textbf{32 of 32 claims verified with no attacks} under
a Dolev--Yao adversary in 0.19\,s. The final system authenticates a hop in 0.4\,ms of
computation at 48.8\,\uJ{} per cycle, sustains authentication across a lossy channel, and
reroutes automatically when a buoy fails.

\vspace{4mm}\noindent
\textbf{Keywords:} underwater acoustic networks, mutual authentication, elliptic curve
cryptography, ECDH, AES-GCM, Scyther, BAN logic, Thorp absorption, replay resistance.

%----------------------------------------------------------------------
\chapter*{Acknowledgement}
\addcontentsline{toc}{chapter}{Acknowledgement}
\thispagestyle{plain}

\noindent
We thank \textbf{Prof.\ Mahendra Shukla}, Department of Information and Communication
Technology, ABV-IIITM Gwalior, for supervising this work. His insistence that we justify
every number in the simulation --- rather than accept a plausible-looking constant --- is
directly responsible for the channel model in Chapter~6 being physical instead of arbitrary.

\vspace{2mm}\noindent
We are grateful to the authors of the base paper for a clearly specified protocol, and to
Prof.\ Cas Cremers and the Scyther project for making rigorous protocol verification
accessible to undergraduates. Finally, we thank the Institute for access to the IEEE Xplore
digital library, without which the base paper would have been out of reach.

%----------------------------------------------------------------------
\cleardoublepage
\pdfbookmark[0]{Contents}{toc}
\tableofcontents
\listoffigures
\addcontentsline{toc}{chapter}{List of Figures}
\listoftables
\addcontentsline{toc}{chapter}{List of Tables}
\cleardoublepage
